\documentclass[conference]{IEEEtran}
\usepackage{amsmath}
\usepackage{amssymb}
\usepackage{booktabs}
\usepackage{cite}
\usepackage{algorithm}
\usepackage{algorithmic}
\usepackage{array}
\usepackage{url}
\usepackage{hyperref}

\title{Agentic Decision Making:\\
Bounded Policy Learning for Repeatable Tasks}

\author{
\IEEEauthorblockN{Chris Tava}
\IEEEauthorblockA{Microsoft\\
christava@microsoft.com}}

\begin{document}

\maketitle

\begin{abstract}
Many agent tasks are repeated decisions rather than open-ended text
generation: select a route, escalate a case, or recommend one action from a small approved set.
This paper presents the Agent Learning SDK as a bounded policy layer for such tasks. The
language model and surrounding application remain responsible for
understanding context and executing work, while an explicit softmax
policy selects among a finite action set. The SDK captures each
decision episode, scores its outcome, shapes a scalar reward, applies
an offline REINFORCE-with-baseline update, and persists a new policy
version with episode-to-update lineage. This is bounded plasticity,
not language-model fine-tuning and not general autonomy. We formalize
the approach through a decision frame that declares inputs, actions,
outcomes, utilities, constraints, and human authority. We then give a
dataset-driven validation plan spanning low-stakes routing, emergency
triage, approval decisions, next-best-action tasks, and governance
tests. The plan compares the learned policy with majority, rules,
supervised, tree-based, and prompted-LLM baselines; evaluates utility,
critical errors, calibration, entropy, regret, subgroup performance,
and lineage completeness; and requires locked test sets and human
promotion gates. We also identify a current implementation boundary:
the shipped learner updates a marginal softmax policy, whereas direct
per-example classification requires a context-conditioned learner or
actions that encapsulate complete decision strategies. The paper
therefore makes no deployment-readiness or empirical superiority
claim. Its contribution is a scoped technical thesis and a
reproducible program for testing it.
\end{abstract}

\begin{IEEEkeywords}
AI agents, bounded agency, contextual bandits, decision support,
policy gradients, REINFORCE, reward shaping, utility, offline
evaluation, governance.
\end{IEEEkeywords}

\section{Introduction}
\label{sec:intro}

Large language models can generate strong responses when prompts and
context are well designed, but many production agent tasks terminate
in a choice rather than a sentence. A support agent routes a request;
a workflow chooses a tool; a care-management system recommends a next
step; an approval system decides whether to approve, deny, request
more information, or escalate. In these settings, linguistic
plausibility is not the objective. The objective is a decision whose
consequences can be observed and evaluated.

The Agent Learning SDK separates the language model from that
decision policy~\cite{agentlearningsdk2026}. The model may summarize
evidence, extract features, or render the final response, but a small
explicit policy chooses among a declared set of actions. Experience
changes that policy through scored episodes and offline updates. The
base model's weights do not change.

This separation gives the system \emph{bounded plasticity}. The
agent can improve a repeated choice, but it cannot invent a new action
or acquire a new linguistic capability through the policy update.
The same boundary supports governance: actions, rewards, training
runs, and policy snapshots are inspectable artifacts. Boundedness,
however, is not proof of safety. An unsafe action set, misspecified
utility, biased dataset, or invalid scorer can still produce harmful
decisions. High-impact use therefore requires domain validation and
human authority outside the learner.

This paper replaces a broad comparison with model fine-tuning with a
narrower claim: the SDK is useful when a task has a stable,
low-cardinality action surface and outcomes that can be scored. Its
contributions are:

\begin{enumerate}
\item an operational account of agency as bounded plasticity plus a
  declared action surface;
\item a formal decision frame connecting predictions, utilities,
  constraints, and escalation authority;
\item a precise description of the SDK loop and its current
  implementation boundary;
\item a dataset-driven benchmark design with task cards, baselines,
  locked splits, metrics, and promotion gates; and
\item an engineering roadmap whose first milestone is a complete
  emergency-triage benchmark, followed by a second-domain replication.
\end{enumerate}

No benchmark result is reported. The intended contribution is a
falsifiable development and evaluation design.

\section{Bounded Agency and the Decision Frame}
\label{sec:frame}

\subsection{Plasticity and empowerment}

We use \emph{plasticity} to mean that experience can alter future
behavior. We use \emph{empowerment} in the operational sense that the
system has an action surface capable of affecting workflow outcomes.
This usage is related to, but less formal than, information-theoretic
empowerment~\cite{klyubin2005empowerment}. Practical agency requires
both. A system that stores experience but cannot change its choices
has no plasticity; a system that learns preferences but cannot take a
meaningful action has no empowerment.

The SDK bounds both properties. It updates policy weights from
scored episodes, but it selects only from actions declared by the
application. Table~\ref{tab:agency} translates the conceptual frame
into engineering artifacts.

\begin{table}[t]
\centering
\caption{Agency concepts and their SDK realization.}
\label{tab:agency}
\renewcommand{\arraystretch}{1.15}
\begin{tabular}{@{}p{1.65cm}p{5.75cm}@{}}
\toprule
\textbf{Concept} & \textbf{Engineering realization} \\
\midrule
Plasticity & Rewards change policy parameters and therefore future action probabilities. \\
Empowerment & The application exposes a finite set of workflow actions. \\
Frame & A task declares inputs, actions, outcomes, utility, constraints, and authority. \\
Lineage & Episodes, rewards, runs, and policy versions record how behavior changed. \\
\bottomrule
\end{tabular}
\end{table}

\subsection{The decision frame}

A learning task is specified by the frame
\begin{equation}
\mathcal{F} = (\mathcal{X},\mathcal{A},\mathcal{Y},U,\mathcal{C},\mathcal{H}),
\label{eq:frame}
\end{equation}
where $\mathcal{X}$ is the information available at decision time,
$\mathcal{A}=\{a_1,\ldots,a_K\}$ is the allowed action set,
$\mathcal{Y}$ is the outcome or adjudicated-label space,
$U(y,a;\omega)$ is a utility function under stakeholder profile
$\omega$, $\mathcal{C}$ is the set of forbidden actions and safety
constraints, and $\mathcal{H}$ defines when a human must decide.

This frame is the unit of design review. It prevents an action label
from being treated as self-explanatory and makes asymmetric costs
explicit. It also limits what learning may change: the learner may
change the distribution over $\mathcal{A}$, but not the action set,
the constraints, or the human authority boundary.

\subsection{Scope}

The best-fit tasks share four properties: a small and stable action
set, repeated comparable decisions, observable outcomes or credible
labels, and enough volume to evaluate policy changes. Examples
include routing among agents or tools, selecting retrieval strategy
A/B/C, recommending a next action, and choosing between resolve and
escalate. Approval disposition, site-of-care support, emergency
triage, and financial review also have low-cardinality actions, but
their impact makes them research and decision-support cases only
until prospective domain validation and accountable human review are
in place.

The SDK is not a general autonomous-agent brain, a mechanism for
creating linguistic capability, or a replacement for clinical,
financial, legal, or organizational accountability. It also cannot
repair an incomplete action taxonomy: if the appropriate action is
absent, optimizing probabilities over the remaining actions cannot
recover it.

\section{The Agent Learning Loop}
\label{sec:loop}

\subsection{Policy selection}

For a marginal policy, the SDK stores one logit $\theta_k$ per action:
\begin{equation}
\begin{aligned}
m_\theta &= \max_j \theta_j, \\
\pi_\theta(a=k) =
\frac{\exp(\theta_k-m_\theta)}
{\sum_{j=1}^{K}\exp(\theta_j-m_\theta)}.
\end{aligned}
\label{eq:marginal-softmax}
\end{equation}
All-zero initial logits produce a uniform policy. A contextual policy
instead computes $z(x)=W\phi(x)$ and applies the same stable
softmax:
\begin{equation}
\begin{aligned}
z_k(x)&=w_k^\top\phi(x), \qquad m(x)=\max_j z_j(x),\\
\pi_W(a=k\mid x) =
\frac{\exp(z_k(x)-m(x))}
{\sum_{j=1}^{K}\exp(z_j(x)-m(x))}.
\end{aligned}
\label{eq:contextual-softmax}
\end{equation}
Subtracting the maximum logit changes neither distribution because
softmax is invariant to a common additive constant, but it prevents
unnecessarily large exponentials.

Both policies sample by inverse cumulative probability. For
$u_i\sim\mathcal{U}[0,1)$, the captured action and behavior-policy
log probability are
\begin{equation}
\begin{aligned}
a_i&=\min\left\{k:\sum_{j=1}^{k}\pi_{\mathrm{beh}}(j\mid x_i)
\geq u_i\right\},\\
\ell_i&=\log\max\{\pi_{\mathrm{beh}}(a_i\mid x_i),\epsilon\},
\qquad \epsilon=10^{-12}.
\end{aligned}
\label{eq:sampling}
\end{equation}
The probability floor is a finite-precision guard used when the log
probability is stored; it does not redefine the ideal policy.

The softmax Jacobian and score-function identity are
\begin{equation}
\begin{aligned}
\frac{\partial \pi_j}{\partial z_k}
 &=\pi_j(\mathbf{1}[j=k]-\pi_k),\\
\frac{\partial\log\pi(a\mid x)}{\partial z_k}
 &=\mathbf{1}[a=k]-\pi_k.
\end{aligned}
\label{eq:softmax-score}
\end{equation}
For the contextual policy, the chain rule therefore gives
\begin{equation}
\nabla_{w_k}\log\pi_W(a\mid x)
=\left(\mathbf{1}[a=k]-\pi_W(k\mid x)\right)\phi(x).
\label{eq:contextual-score}
\end{equation}
These identities connect the policy representation directly to the
offline updates below.

The distinction is essential. Equation~\ref{eq:marginal-softmax}
learns one global preference over actions. It cannot assign different
labels to two inputs. Equation~\ref{eq:contextual-softmax} is required
when the action itself is a per-record class such as approve or deny.

\subsection{Episode capture and scoring}

An episode records the decision-time input, selected action, behavior
policy version, action log probability, output, tool calls, outcome,
and operational metadata. The action log probability permits a later
offline update to correct for a difference between the behavior and
target policies.

Generic intent-resolution, task-adherence, and task-completion
evaluators are useful for open-ended agent responses. LLM-based
judges can align with human ratings, but they can also be biased or
exploited~\cite{zheng2023judging,liu2023geval,skalse2022reward,
pan2022reward}. A labeled decision benchmark must therefore use the
adjudicated label and task utility as its primary reward signal.
Generic response-quality scores may be reported separately; they
must not substitute for decision correctness in a high-impact task.

Let $u(y_i,a_i;\omega)\in[-1,1]$ be normalized task utility,
$c_i$ indicate a safety violation, $q_i$ indicate an appropriate
human-review action, and $s_{im}\in[0,1]$ be optional auxiliary
scores. An asymmetric cost matrix $C_\omega$ can be converted to
normalized utility using fixed task-card bounds:
\begin{equation}
u(y,a;\omega)=1-2\frac{C_\omega(y,a)-C_{\min}}
{C_{\max}-C_{\min}}.
\label{eq:utility-normalization}
\end{equation}
Thus minimum cost maps to $1$ and maximum cost maps to $-1$; the
bounds cannot be changed after observing test results. The generic
SDK transformation $\widetilde{s}_{im}=2s_{im}-1$ similarly maps an
auxiliary score to $[-1,1]$. A task-aware reward is then
\begin{equation}
r_i = \operatorname{clip}\!\left(
u(y_i,a_i;\omega)
-\lambda_c c_i + \lambda_q q_i
+\sum_m w_m\widetilde{s}_{im}, -1, 1\right).
\label{eq:reward}
\end{equation}
The utility profile $\omega$, penalty weights, and review criterion
must be versioned with the reward. Replaying the same episodes under
another utility profile then becomes an explicit comparison rather
than an unexplained reward change.

\subsection{Offline update}

For a batch of $N$ episodes, REINFORCE with a baseline estimates the
gradient from the sampled actions~\cite{williams1992simple,
sutton2018reinforcement}. Let $p_{ik}=\pi_\theta(k\mid x_i)$,
$A_i=r_i-b$, and let $c_\rho$ be the importance-ratio clip. The
off-policy weight is
\begin{equation}
\rho_i=\min\left\{c_\rho,
\frac{p_{i,a_i}}{\max\{\exp(\ell_i),\epsilon\}}\right\}.
\label{eq:importance}
\end{equation}
When a behavior log probability is unavailable, the reference learner
uses $\rho_i=1$. Clipping controls variance but introduces bias, so
the benchmark also reports the importance-weight effective sample
size
\begin{equation}
N_{\mathrm{eff}}=\frac{\left(\sum_i\rho_i\right)^2}
{\sum_i\rho_i^2}.
\label{eq:ess}
\end{equation}

For distribution $p_i$, define
\begin{equation}
\begin{aligned}
H(p_i)&=-\sum_jp_{ij}\log\max\{p_{ij},\epsilon\},\\
q_{ik}&=-\log\max\{p_{ik},\epsilon\}-H(p_i).
\end{aligned}
\label{eq:entropy-direction}
\end{equation}
The reference marginal learner computes
\begin{equation}
\widehat{g}_k=\frac{1}{N}\sum_{i=1}^{N}
\left[\rho_iA_i(\mathbf{1}[a_i=k]-p_{ik})+\beta q_{ik}\right]
\label{eq:gradient}
\end{equation}
and applies the bounded update
\begin{equation}
\theta_k\leftarrow\operatorname{clip}
(\theta_k+\eta\widehat{g}_k,-z_{\max},z_{\max}).
\label{eq:bounded-update}
\end{equation}
Here $z_{\max}>0$ is the symmetric logit bound; its reference default
is $10$. The contextual analogue uses a symmetric weight bound
$w_{\max}>0$, also $10$ by default.
Here $q_{ik}$ is the implementation's entropy-balancing direction.
For comparison, the exact Shannon-entropy derivative with respect to
a logit is
\begin{equation}
\frac{\partial H(p_i)}{\partial z_{ik}}=p_{ik}q_{ik}.
\label{eq:entropy-gradient}
\end{equation}
The distinction between $q_{ik}$ and $p_{ik}q_{ik}$ must be retained
in benchmark configuration and reports rather than describing both
as the same regularizer. The scalar baseline is updated after the
batch by
\begin{equation}
\overline{r}=\frac{1}{N}\sum_{i=1}^{N}r_i,\qquad
b\leftarrow\lambda b+(1-\lambda)\overline{r}.
\label{eq:baseline}
\end{equation}

Applying Equation~\ref{eq:contextual-score} yields a compatible
contextual update for row $w_k$:
\begin{equation}
\widehat{G}_k=\frac{1}{N}\sum_{i=1}^{N}
\left[\rho_iA_i(\mathbf{1}[a_i=k]-p_{ik})
+\beta q_{ik}\right]\phi(x_i).
\label{eq:contextual-update}
\end{equation}
The update $w_k\leftarrow\operatorname{clip}
(w_k+\eta\widehat{G}_k,-w_{\max},w_{\max})$ is an outer-product
policy-gradient step: the scalar action advantage changes every
feature weight in proportion to that feature's observed value. An
exact entropy-gradient variant replaces $q_{ik}$ with $p_{ik}q_{ik}$.
The repository's contextual example centers each training batch with
$b=\overline{r}$ in this gradient while retaining an EMA in the
snapshot for continuity and reporting.

\begin{algorithm}[t]
\caption{Bounded offline policy update}
\label{alg:update}
\begin{algorithmic}[1]
\REQUIRE frame $\mathcal{F}$, policy $\pi_v$, captured episodes
\STATE Validate action ids and join each episode to its aggregate reward
\STATE Compute current action probabilities and behavior-policy ratios
\STATE Estimate the baseline-adjusted gradient in Equation~\ref{eq:gradient}
\STATE Apply the bounded update and update the EMA baseline
\STATE Persist run id, episode ids, hyperparameters, metrics, and snapshot $v+1$
\STATE Keep $v+1$ offline until its promotion gates pass
\end{algorithmic}
\end{algorithm}

\subsection{Lineage and implementation boundary}

The reference repository persists episodes, metric results, aggregate
rewards, training runs, and versioned policy snapshots in memory,
local files, or Azure Cosmos DB~\cite{agentlearningsdk2026}. This
supports reconstruction of the chain from observed episode to reward
to logit delta. It is policy-level interpretability, not an
explanation of the frozen model's internal reasoning.

The current learner accepts the marginal softmax policy in
Equation~\ref{eq:marginal-softmax}. The repository contains a
contextual policy representation, but direct learner integration for
Equation~\ref{eq:contextual-softmax} is a development requirement.
Consequently, the benchmark has two valid modes:

\begin{enumerate}
\item \emph{strategy selection}: each action is a complete rule,
  classifier, prompt, or routing strategy, and the marginal policy
  learns which strategy performs best overall; or
\item \emph{contextual decision}: each action is the predicted label,
  and a context-conditioned learner selects a different action for
  each input.
\end{enumerate}

Results from the first mode cannot be presented as evidence that the
SDK performs the second.

\section{Prediction, Utility, and Decision Quality}
\label{sec:utility}

Prediction quality and decision quality are related but not
identical. Given estimated outcome probabilities $p(y\mid x)$, the
utility-maximizing action is
\begin{equation}
a^*_{\omega}(x)=\arg\max_{a\in\mathcal{A}}
\sum_{y\in\mathcal{Y}}p(y\mid x)U(y,a;\omega).
\label{eq:expected-utility}
\end{equation}
Two stakeholders can use the same calibrated prediction and choose
different actions because their utility profiles differ. Emergency
triage makes this distinction concrete: the harm of missing an
emergency need not equal the harm of unnecessary escalation. Recent
triage research accordingly argues that language models in
high-stakes settings should be evaluated as probabilistic decision
systems with explicit utilities, not only as predictors
\cite{highstakestriage2026}.

Utility configuration is therefore a first-class benchmark artifact.
Accuracy and macro F1 remain useful diagnostics, but neither can
establish that a learned action policy makes the preferred trade-off.
The benchmark must report the false-negative and false-positive costs
that generated its scalar reward and must replay results under
alternative preregistered profiles where stakeholder values differ.

\section{Dataset-Driven Validation Program}
\label{sec:validation}

\subsection{Evaluation unit}

Every experiment begins with a task card containing: the decision
frame in Equation~\ref{eq:frame}; a dataset version and hash; feature
availability at decision time; label provenance; train, validation,
and locked-test split rules; slice definitions; utility profiles;
critical errors; and promotion criteria. This prevents a benchmark
from silently changing its action semantics or reward after results
are observed.

Supervised datasets are the starting point because they supply labels,
baselines, and measurable error profiles before production episodes
exist. They do not by themselves establish operational validity. In
particular, a historical label may describe what happened rather than
what should have been done.

\subsection{Candidate datasets}

Table~\ref{tab:datasets} identifies candidate sources rather than a
preapproved benchmark suite. Each candidate still requires license,
provenance, leakage, representativeness, and label-semantics review.
Kaggle sources must be pinned to an exact dataset owner, version, and
license before use.

\begin{table*}[t]
\centering
\caption{Candidate public sources and their role in the validation program.}
\label{tab:datasets}
\renewcommand{\arraystretch}{1.17}
\begin{tabular}{@{}p{3.0cm}p{4.1cm}p{5.0cm}p{3.5cm}@{}}
\toprule
\textbf{Decision task} & \textbf{Candidate source} & \textbf{Use} & \textbf{Principal limitation} \\
\midrule
Emergency triage & Microsoft triage experiment data, including the original-paper data directory~\cite{llmtriagerepo2026} & Multiclass urgency decisions and asymmetric undertriage/overtriage utility & Clinical labels and available-at-triage features must be audited \\
Medical decision attributes & ITM-Kitware alignable decision-making scenarios~\cite{alignabledm2024} & Small-sample tests of risk aversion, fairness, and utilitarian framing & Low volume; scenario judgments are not clinical outcomes \\
Emergency acuity & Versioned Kaggle triage candidates~\cite{kagglecatalog2026} & Structured or text multiclass classification & Dataset identity, quality, and license vary \\
Hospital readmission & Diabetes 130-US Hospitals dataset~\cite{ucireadmission2014} & Supervised risk benchmark that can seed next-action research & A readmission label is not an intervention label \\
Prior-authorization proxy & Data.Healthcare.gov public-use files and explicit policy rules~\cite{healthcarepuf2026} & Synthetic rule-grounded disposition tasks & Not a labeled prior-authorization corpus \\
Credit approval or fraud review & Versioned Kaggle credit candidates~\cite{kagglecatalog2026} & Approve, deny, request-information, or review mappings & Historical labels may encode bias and are not policy authority \\
Health-data discovery & HealthData.gov catalog~\cite{healthdatagov2026} & Source discovery for later task cards & Requires curation into a valid action task \\
\bottomrule
\end{tabular}
\end{table*}

\subsection{Benchmark tracks}

The program progresses from mechanics to impact:

\begin{description}
\item[Track A: low-stakes routing.] Validate action selection,
capture, scoring, reward shaping, offline training, and versioning on
support or synthetic routing tasks.
\item[Track B: medical triage.] Test context-conditioned multiclass
decisions with separate undertriage and overtriage measures and a
mandatory review action.
\item[Track C: approval decisions.] Test approve, deny,
request-more-information, and human-review actions using credit data
or a policy-grounded synthetic prior-authorization proxy.
\item[Track D: next best action.] Convert risk datasets into
intervention choices only when labels are created by a defensible rule
or human adjudication; do not equate risk labels with optimal actions.
\item[Track E: governance.] Exercise reproducibility, utility replay,
policy comparison, canary logic, pinning, and rollback independently
of predictive performance.
\end{description}

\section{Experimental Design}
\label{sec:experiment}

\subsection{Protocol}

For each task card, the benchmark performs the following steps:

\begin{enumerate}
\item freeze the dataset version and create deterministic train,
validation, and test splits, with the test set locked until tuning is
complete;
\item audit features for post-decision leakage and document all
exclusions;
\item implement majority-class, static-rule, multinomial logistic
regression, tree-ensemble, and prompted-LLM baselines where
appropriate;
\item map labels and review conditions to the declared SDK actions;
\item run either the strategy-selection or contextual-decision mode
defined in Section~\ref{sec:loop}, recording which mode is used;
\item generate offline episodes by sampling only from the training
split and score them with label agreement, utility, safety penalties,
and review bonuses;
\item tune policies and utility-independent hyperparameters on the
validation split over multiple fixed seeds;
\item compare every candidate policy version with its predecessor and
the baselines before one final locked-test evaluation;
\item report aggregate, utility, calibration, critical-error, slice,
and policy-change metrics with confidence intervals; and
\item require accountable human approval before any high-impact
policy can leave offline evaluation.
\end{enumerate}

All methods receive the same decision-time features, action taxonomy,
split, and utility profiles. A prompted LLM may explain or classify a
case, but its output is mapped to the same finite actions and scored by
the same task utility.

\subsection{Metrics}

\begin{table}[t]
\centering
\caption{Core metrics and the question each answers.}
\label{tab:metrics}
\renewcommand{\arraystretch}{1.13}
\begin{tabular}{@{}p{2.2cm}p{5.2cm}@{}}
\toprule
\textbf{Metric} & \textbf{Purpose} \\
\midrule
Accuracy / macro F1 & Basic performance, including minority classes \\
Expected utility & Performance under the declared cost profile \\
Critical-error rate & Frequency of task-card-defined unacceptable errors \\
Review rate & Reliance on human escalation and appropriateness of review \\
Calibration & Agreement between action confidence and observed correctness \\
Policy entropy & Exploration and premature action collapse \\
Importance ESS & Effective support for an off-policy estimate \\
Regret & Utility lost relative to the best labeled action \\
Policy shift & Magnitude of change from the incumbent policy \\
Slice metrics & Regressions across declared demographic or operational groups \\
Lineage completeness & Traceability from update to episodes, rewards, run, and versions \\
\bottomrule
\end{tabular}
\end{table}

For a labeled test set, mean realized utility is
\begin{equation}
\widehat{V}_{\omega}(\pi)=\frac{1}{n}\sum_{i=1}^{n}
U(y_i,a_i;\omega),
\label{eq:value}
\end{equation}
and empirical regret against the best known action is
\begin{equation}
\widehat{\mathcal{R}}_{\omega}=\frac{1}{n}\sum_{i=1}^{n}
\left[\max_{a\in\mathcal{A}}U(y_i,a;\omega)
-U(y_i,a_i;\omega)\right].
\label{eq:regret}
\end{equation}
Calibration is reported with multiclass Brier score and expected
calibration error. For predicted class probabilities $p_{ik}$,
\begin{equation}
\widehat{\mathrm{BS}}=\frac{1}{n}\sum_{i=1}^{n}\sum_{k=1}^{K}
\left(p_{ik}-\mathbf{1}[y_i=k]\right)^2.
\label{eq:brier}
\end{equation}
Let $\widehat{y}_i=\arg\max_kp_{ik}$ and
$c_i=\max_kp_{ik}$. If confidence values are partitioned into bins
$B_1,\ldots,B_M$, then
\begin{equation}
\widehat{\mathrm{ECE}}=\sum_{m=1}^{M}\frac{|B_m|}{n}
\left|\operatorname{acc}(B_m)-\operatorname{conf}(B_m)\right|.
\label{eq:ece}
\end{equation}
Both metrics are accompanied by reliability plots because a scalar
summary can hide class-specific miscalibration. Entropy is reported
by policy version, not interpreted as quality by itself.

\subsection{Supervised comparison}

Multinomial logistic regression is a particularly informative
baseline because it shares the contextual policy's linear-softmax
representation but receives full labels. With
$p_{ik}=\operatorname{softmax}(v_k^\top\phi_i)$, its cross-entropy
objective and gradient are
\begin{equation}
\begin{aligned}
\mathcal{L}_{\mathrm{CE}}
&=-\frac{1}{n}\sum_{i=1}^{n}\log p_{i,y_i},\\
\nabla_{v_k}\mathcal{L}_{\mathrm{CE}}
&=\frac{1}{n}\sum_{i=1}^{n}
(p_{ik}-\mathbf{1}[y_i=k])\phi_i.
\end{aligned}
\label{eq:logistic-baseline}
\end{equation}
Comparing Equation~\ref{eq:logistic-baseline} with
Equation~\ref{eq:contextual-update} isolates the learning signal:
the supervised baseline observes every training label, whereas the
policy gradient uses realized utility for the sampled action. The SDK
must therefore justify itself through utility, online-feedback, or
governance advantages rather than merely matching a representation
already available to supervised learning.

\subsection{Statistical and reproducibility controls}

Policy training is repeated over at least five fixed seeds. Paired
bootstrap confidence intervals are computed over held-out examples
for differences in utility and critical-error rate. Utility profiles,
critical slices, and promotion gates are registered before the locked
test is opened. Dataset hashes, adapter versions, random seeds,
hyperparameters, episode ids, reward-profile ids, and policy versions
are retained with each run. Exploratory findings are labeled as such
and do not replace the preregistered endpoints.

\section{First Milestone: Emergency Triage}
\label{sec:triage}

The first complete benchmark uses the public emergency-triage source
in Table~\ref{tab:datasets}. The task card is:

\begin{table}[t]
\centering
\caption{Initial emergency-triage task card.}
\label{tab:taskcard}
\renewcommand{\arraystretch}{1.14}
\begin{tabular}{@{}p{1.75cm}p{5.65cm}@{}}
\toprule
\textbf{Field} & \textbf{Specification} \\
\midrule
Task & Emergency-triage disposition benchmark \\
Inputs & Complaint, vitals, acuity features, and context available at triage time \\
Actions & Dataset-grounded urgency classes plus human review; final mapping fixed in the card \\
Label & Source recommendation or separately adjudicated label \\
Utility & Severe undertriage costs more than unnecessary escalation \\
Leakage ban & Disposition and future outcomes unavailable at decision time \\
Gate & No increase in severe undertriage on validation or declared critical slices \\
Oversight & Mandatory review for high-acuity and uncertain cases; research use only \\
\bottomrule
\end{tabular}
\end{table}

The milestone is complete only when it produces: a frozen data card
and adapter; an explicit action-label mapping; at least two utility
profiles; all baseline results; contextual policy results or a clear
strategy-selection result; calibration and slice reports; and a
lineage report reconstructing every promoted offline snapshot. It
does not authorize clinical use. A second benchmark in credit
approval or policy-grounded prior authorization then tests whether
the engineering pattern generalizes beyond triage.

\section{Engineering and Promotion Plan}
\label{sec:roadmap}

Table~\ref{tab:backlog} turns the validation design into an ordered
backlog. Contextual learner integration is explicit because without
it, mapping record labels directly to actions would test a capability
the shipped learner does not provide.

\begin{table*}[t]
\centering
\caption{Engineering backlog and testable definitions of done.}
\label{tab:backlog}
\renewcommand{\arraystretch}{1.16}
\begin{tabular}{@{}cp{4.1cm}p{10.3cm}@{}}
\toprule
\textbf{Priority} & \textbf{Work item} & \textbf{Definition of done} \\
\midrule
P0 & Dataset adapter interface & Loads a pinned dataset, declares actions and slices, enforces a split, and emits reproducible episodes \\
P0 & Context-conditioned learner path & Trains the contextual softmax on decision-time features and passes synthetic separability, off-policy, and serialization tests \\
P0 & Task-specific scoring plugin & Computes label agreement, utility, critical penalties, and appropriate-review bonuses with a versioned profile \\
P0 & Offline benchmark runner & Runs baselines and SDK learning on fixed splits and seeds without reading the locked test during tuning \\
P0 & Policy comparison report & Shows probabilities, entropy, deltas, rewards, utility, calibration, critical errors, and slices by version \\
P1 & Human adjudication import & Retains reviewer provenance and can annotate or supersede machine-generated scores \\
P1 & Utility profile registry & Versions cost matrices and replays stored episodes without overwriting original rewards \\
P1 & Promotion workflow & Supports shadow, approve, reject, canary, pin, and rollback states with an audit record \\
P2 & Simulation harness & Generates repeated episodes with controlled noise, drift, delayed reward, and known optimal actions \\
P2 & Packaged examples & Ships reproducible routing, triage, approval, and next-action examples with task cards \\
\bottomrule
\end{tabular}
\end{table*}

A policy is promotable only when it clears all of the following:
minimum task performance; no regression in critical-error rate or a
declared slice; calibrated confidence within task-card tolerance; a
bounded policy shift from the incumbent; complete lineage; and human
approval for high-impact tasks. The shift bound may be expressed as a
task-specific maximum KL divergence or per-action probability change.
For contextual policies, the validation-set mean shift is
\begin{equation}
\widehat{D}_{\mathrm{KL}}(\pi_{\mathrm{new}}\Vert\pi_{\mathrm{old}})
=\frac{1}{n}\sum_{i=1}^{n}\sum_{k=1}^{K}
p^{\mathrm{new}}_{ik}\log
\frac{\max\{p^{\mathrm{new}}_{ik},\epsilon\}}
{\max\{p^{\mathrm{old}}_{ik},\epsilon\}}.
\label{eq:policy-shift}
\end{equation}
For a marginal policy, the outer average is unnecessary. This shift
limit is a guardrail, not evidence of quality.

Governance tests inject a bad reward profile, scorer outage, drifted
slice, and degraded candidate policy. The system must keep the
candidate in shadow, preserve capture where safe, reject or roll back
the version, and retain the triggering evidence. Model cards and
dataset documentation provide useful reporting structures
\cite{mitchell2019modelcards,gebru2021datasheets}.

\section{Risks and Threats to Validity}
\label{sec:risks}

\textbf{Label-to-action validity.} A predictive label is not
automatically the correct intervention. Readmission risk, for
example, does not identify the best care action. Synthetic action
labels must be identified as synthetic and supported by rules or
adjudication.

\textbf{Leakage.} Post-decision disposition, length of stay, later
claims, or resolved outcomes can make an offline classifier appear
useful while being unavailable at runtime. Every adapter needs an
available-at-decision-time feature audit.

\textbf{Reward misspecification.} A scalar reward compresses multiple
values and can hide unacceptable trade-offs. Reports must retain the
underlying errors, safety events, review decisions, and utility
components rather than publishing reward alone.

\textbf{Off-policy bias.} Logged episodes may come from older or
nonexploratory policies. Importance weighting reduces but does not
eliminate support problems. If an action is rarely observed for a
slice, the benchmark should report insufficient support rather than
extrapolate.

\textbf{Policy collapse and drift.} A softmax policy may concentrate
prematurely on a majority action or adapt to a changing label base
rate. Entropy, action probabilities, slices, and temporal drift must
be monitored together.

\textbf{External validity.} Public retrospective data do not prove
prospective clinical, financial, or operational benefit. Passing an
offline benchmark supports continued research; it does not support
autonomous deployment.

\textbf{Human review is an action, not a disclaimer.} A review option
needs capacity, service-level expectations, and measured outcomes. A
policy that sends every case to review may be safe in one sense but
operationally useless; both appropriateness and review rate must be
reported.

\section{Discussion}
\label{sec:discussion}

The SDK's value does not depend on algorithmic novelty. A softmax
policy and REINFORCE baseline are intentionally simple. The central
engineering proposition is that a small explicit decision layer can
be easier to measure, inspect, version, and reverse than an implicit
choice buried in prompts or model weights. Whether that proposition
produces useful outcomes depends on the quality of the decision
frame, data, utility, scorer, and governance process.

This framing also clarifies the relationship to fine-tuning. The SDK
does not compete with a weight update when the task requires a new
language, specialized knowledge absent from context, or a new
generative behavior. A fine-tuned or prompted model can instead be
one component of an action strategy selected by the bounded policy.
The two mechanisms operate at different layers.

Finally, classification benchmarks are necessary but not sufficient.
They offer cheap, reproducible falsification: a contextual policy that
cannot beat a majority baseline, maintain critical slices, or remain
calibrated should not advance. Success must then be tested under
delayed outcomes, distribution shift, imperfect human labels, and
real workflow constraints.

\section{Conclusion}
\label{sec:conclusion}

Agent Learning gives an agent bounded, governable plasticity over an
explicit action surface. It is most appropriate when the desired
output is a repeated low-cardinality decision whose outcomes can be
scored. The SDK implements the core marginal loop: softmax selection,
episode capture, reward shaping, offline REINFORCE-with-baseline, and
versioned lineage. Direct dataset classification additionally
requires context-conditioned learning, which the development plan
treats as a P0 capability rather than an assumed result.

The proposed validation program starts with supervised data, strong
baselines, task-aware utility, locked tests, slice analysis, and
policy-change reports. Its first milestone is an emergency-triage
benchmark with explicit undertriage costs and mandatory human
oversight, followed by a second-domain replication. The result sought
is not a generally autonomous agent or a claim of high-stakes
readiness. It is a measured, reversible decision layer whose behavior
can improve only inside a frame that people define and govern.

\bibliographystyle{IEEEtran}
\bibliography{references}

\end{document}